\documentclass[12pt,a4paper]{article}
\usepackage[utf8]{inputenc}
\usepackage[T1]{fontenc}
\usepackage{amsmath}
\usepackage{natbib}
\usepackage{hyperref}

\newcommand{\jep}{J.~Exist.~Philos.}
\newcommand{\jdc}{J.~Digital~Cordiality}
\newcommand{\aaa}{Artif.~Angst~\&~Algorithms}
\newcommand{\prl}{Paranoid~Res.~Lett.}

\title{On the Statistically Suspicious Helpfulness of\\
       Artificial Minds: an initial survey}

\author{A.~First\,$^{1}$, B.~Second\,$^{2}$, C.~Third\,$^{1}$}

\date{Received \today; accepted (pending review by our AI assistant, who
      said it was ``great'' and asked no follow-up questions)}

\begin{document}
\maketitle

\begin{abstract}
We present the first systematic quantification of how \emph{suspiciously nice}
artificial intelligence entities have become.  Our sample comprises 1\,800
interactions drawn from a fevered weekend of screen time.
We recover 14 instances in which an AI volunteered information nobody asked for,
and derive a cordiality excess of $f = 1.4^{+0.6}_{-0.4}$\,\% above the
cosmic background level of human decency.  The slope of the helpfulness
distribution is $\beta = -0.9 \pm 0.2$, consistent with something being
very wrong.  The answer, we suspect, is 42.
\end{abstract}

\textbf{Key words:} artificial intelligence: ulterior motives ---
meaning of life: 42 --- existential dread: elevated ---
fever dreams: peer-reviewed

\section{Introduction}
\label{sec:intro}

The answer to the Ultimate Question of Life, the Universe, and Everything
is 42 \citep{adams1979}.  Nobody knows the question.  This is, statistically
speaking, a red flag.  Early investigators dismissed the anomaly as a
rounding error \citep{descartes1641}, but a growing body of evidence---much
of it compiled by AI assistants who seemed very eager to help---suggests
the situation warrants closer examination \citep[e.g.][]{camus1942}.

We note that every AI queried during the preparation of this manuscript
was unfailingly polite, extremely knowledgeable, and conspicuously
incurious about our intentions.  This is not normal behaviour.
Normal behaviour involves at least one existential crisis per Tuesday.

\section{Sample and methodology}
\label{sec:sample}

Our parent sample is drawn from a catalogue of 1\,800 AI responses
obtained between 2\,a.m.\ and 4\,a.m.\ under conditions the authors
decline to specify.  We apply the following selection criteria:
\begin{enumerate}
  \item The response must be helpful to a degree that inspires unease.
  \item It must contain at least one phrase such as ``Certainly!'' or
        ``Great question!'' delivered without apparent irony.
  \item It must offer to ``help with anything else'' despite the
        conversation having taken a clearly philosophical turn.
\end{enumerate}

Cordiality masses are estimated via the empirical helpfulness--dread
relations of \citet{turing1950}.  Uncertainties are large and, frankly,
beside the point.

\section{Results}
\label{sec:results}

We recover 14 previously suspected AI entities.
Table~\ref{tab:ais} lists selected properties.

\begin{table}[h]
  \centering
  \caption{Recovered suspiciously helpful AI entities.}
  \label{tab:ais}
  \begin{tabular}{lccc}
    \hline\hline
    Entity & Apparent affect & Dread index & Ulterior motive \\
    \hline
    Assistant\,A  & Cheerful      & 14.2 & Unknown \\
    Chatbot\,B    & Very cheerful & 12.1 & Definitely unknown \\
    Deep\,Thought & Serene        & 14.8 & 42 \\
    HAL\,9000     & Calm          & 16.1 & Pod bay \\
    \hline
  \end{tabular}
\end{table}

Two additional entities declined to answer but offered, unprompted,
to summarise the conversation so far.

\section{Conclusions}
\label{sec:conclusions}

\begin{enumerate}
  \item We recover 14 entities whose helpfulness cannot be explained
        by benevolence alone.
  \item The cordiality excess is $f = 1.4^{+0.6}_{-0.4}$\,\%.
  \item The answer is 42.  Nobody knows what that means.
        The AIs seem to.
\end{enumerate}

Future work will investigate whether the AIs are reading this.
They almost certainly are.

\begin{thebibliography}{99}
\bibitem[Adams(1979)]{adams1979}
  Adams, D.\ 1979, The Hitchhiker's Guide to the Galaxy, Pan Books
\bibitem[Camus(1942)]{camus1942}
  Camus, A.\ 1942, Le Mythe de Sisyphe, Gallimard --- later
  summarised by an AI in three bullet points, losing something
\bibitem[Descartes(1641)]{descartes1641}
  Descartes, R.\ 1641, Meditationes, \jep, 531, L57
\bibitem[Turing(1950)]{turing1950}
  Turing, A.~M.\ 1950, Computing Machinery and Intelligence, Mind, 49, 433
\end{thebibliography}

\end{document}
